\begin{proof}[Proof of \cref{obj:thm:observable-upper}]
\textit{1. Pilot calibration and auxiliary experiment.}
Choose
\[
H_0=1025.
\]
Thus \(H_0>0\), \(H_0\ge1\), and \(H_0\ge1024\), the universal self-normalization multiplier.

We first establish the auxiliary risk bound for integers \(n,d\ge1\), an overlap parameter \(0<\epsilon\le1/2\), and a law \(\mathbb P\in\mathcal D_{d,\epsilon}^{\mathrm{obs}}\). Write
\[
L=L_d=\log(ed),\qquad m=\frac n8,\qquad
\delta_{\mathrm p}=\frac Lm,\qquad
q_x=(q_{ay,x})_{a,y\in\{0,1\}}.
\]
The atom probabilities are those of \cref{obj:def:observed-class}; in particular, \(m>0\) and \(L=1+\log d\ge1\).

Consider the auxiliary experiment that draws \(\mathsf M\sim\operatorname{Pois}(n/4)\) and, conditional on \(\mathsf M\), draws \(\mathsf M\) independent observations with law \(\mathbb P\), together with independent fair marks. Its law on the disjoint union of all finite marked samples is
\[
\mathcal Q_{\mathbb P}
=
\sum_{k=0}^{\infty}
e^{-n/4}\frac{(n/4)^k}{k!}
\left(\mathbb P\otimes\operatorname{Bernoulli}(1/2)\right)^{\otimes k}.
\]
Write \(\mathbb E_{\mathrm{aux}}\) for expectation under this law. Construct the counts and cell statistics of \cref{obj:def:armwise-estimator} for every finite marked sample, with the auxiliary count unrestricted. For arrays
\(\boldsymbol t^{\mathrm p},\boldsymbol t^{\mathrm e}\in[0,1]^{\mathcal J_d}\), indexed by the observed atom set defined in \cref{obj:synth_1}, the count generating function is
\[
\mathbb E_{\mathrm{aux}}
\prod_{x,a,y}
(t^{\mathrm p}_{ay,x})^{N'_{ay,x}}
(t^{\mathrm e}_{ay,x})^{N_{ay,x}}
=
\exp\left\{
m\sum_{x,a,y}q_{ay,x}(t^{\mathrm p}_{ay,x}-1)
+
m\sum_{x,a,y}q_{ay,x}(t^{\mathrm e}_{ay,x}-1)
\right\}.
\]
Indeed, conditional on the auxiliary count, each marked observation contributes one category factor; summing over the Poisson count gives the displayed exponential. Consequently,
\[
N'_{ay,x},N_{ay,x}\sim\operatorname{Pois}(mq_{ay,x}),
\]
and all these counts are mutually independent. In particular, the complete cell statistics \(T_x\) are independent across distinct cells.
% lean: CausalSmith.Stat.OptvalueVanishingoverlapRate.markedCount_poisson_hasLaw
% lean: CausalSmith.Stat.OptvalueVanishingoverlapRate.markedPoissonCellValue_pairwise_indepFun

\textit{2. Polynomial oscillation and coefficient control.}
Use the centers, radii, rectangles, polynomial, and clipping scales of \cref{obj:def:armwise-estimator}. Every coordinate radius is positive, since
\[
h_{ay,x}=H_0\left(\sqrt{c_{ay,x}\delta_{\mathrm p}}+\delta_{\mathrm p}\right)>0,
\qquad
\frac{h_{ay,x}}2\le r_{ay,x}\le h_{ay,x}.
\]
Moreover,
\[
c_{ay,x}\le b_{ay,x}\le c_{ay,x}+h_{ay,x}.
\]
Thus \(b_x\ne0\). Applying \cref{obj:lem:armwise-modulus} with comparison vector \(b_x\), for every \(v\in Q_x\), gives
\[
|F_\epsilon(v)-F_\epsilon(b_x)|
\le
2\sum_{a=0}^1
\left(1+\frac{s(b_x)}{D_a(b_x)}\right)R_{a,x}
=S_x^{\mathrm{aw}}.
\]
The Jackson convolution is positive and normalized, so the same bound holds for its polynomial:
\[
\sup_{v\in Q_x}
|P_{x,K_n}(v)-F_\epsilon(b_x)|
\le S_x^{\mathrm{aw}}.
\]

For completeness, define the Chebyshev polynomials by
\[
\mathsf{Ch}_0(\xi_\circ)=1,\qquad
\mathsf{Ch}_1(\xi_\circ)=\xi_\circ,\qquad
\mathsf{Ch}_{j+1}(\xi_\circ)
=2\xi_\circ\mathsf{Ch}_j(\xi_\circ)-\mathsf{Ch}_{j-1}(\xi_\circ),
\quad j\ge1.
\]
Here \(\xi_\circ\) is a real scalar. The recurrence implies
\[
\sum_{\ell=0}^{j}
\left|[\xi_\circ^\ell]\mathsf{Ch}_j(\xi_\circ)\right|
\le(1+\sqrt2)^j,
\]
by induction and the identity
\((1+\sqrt2)^2=2(1+\sqrt2)+1\).

With \(D=2(K_n-1)\), expand the normalized polynomial as
\[
P_{x,K_n}(b_x+r_x\odot\xi)-F_\epsilon(b_x)
=
\sum_{\boldsymbol\nu\in\{0,\ldots,D\}^4}
\gamma_{x,\boldsymbol\nu}
\prod_{a,y}\mathsf{Ch}_{\nu_{ay}}(\xi_{ay}),
\qquad \xi\in[-1,1]^4.
\]
Cosine orthogonality expresses each coefficient as an integral of the bounded polynomial against cosine factors, with normalization factor at most \(16\). Hence
\[
|\gamma_{x,\boldsymbol\nu}|\le16S_x^{\mathrm{aw}}.
\]
Converting the tensor Chebyshev expansion to the monomial expansion defining \(\beta_{x,\boldsymbol\alpha}\) yields
\[
\sum_{\boldsymbol\alpha}|\beta_{x,\boldsymbol\alpha}|
\le
16(D+1)^4(1+\sqrt2)^{4D}S_x^{\mathrm{aw}}.
\]
The numerical factor satisfies
\[
\left\{16(D+1)^4(1+\sqrt2)^{4D}\right\}^2
\le e^{256+24D}\le e^{304K_n}.
\]
Indeed, \(D+1\le e^D\), \(1+\sqrt2\le e^2\), \(256\le e^{256}\), and \(D\le2K_n\), \(K_n\ge1\).
% lean: CausalSmith.Stat.OptvalueVanishingoverlapRate.jacksonCoeff_sum_abs_le
% lean: CausalSmith.Stat.OptvalueVanishingoverlapRate.jacksonCoeff_factor_square_le_exp

\textit{3. Factorial second moments and universal degree tuning.}
Define the successful-localization event by
\[
\mathcal G_x=\{q_x\in Q_x\}.
\]
On this event,
\[
|q_{ay,x}-b_{ay,x}|\le r_{ay,x},
\qquad
\frac{q_{ay,x}}{m r_{ay,x}^2}\le\frac8L.
\]
To verify the second inequality, the coordinate formulas give
\[
c_{ay,x}\delta_{\mathrm p}\le h_{ay,x}^2,\qquad
\delta_{\mathrm p}\le h_{ay,x},\qquad
r_{ay,x}\ge\frac{h_{ay,x}}2,
\]
because \(H_0\ge1\). Since \(q_{ay,x}\le c_{ay,x}+h_{ay,x}\),
\[
q_{ay,x}\delta_{\mathrm p}
\le(c_{ay,x}+h_{ay,x})\delta_{\mathrm p}
\le2h_{ay,x}^2
\le8r_{ay,x}^2,
\]
which proves the claim.

For a Poisson count \(\mathsf N_\circ\) with mean \(m q_\circ\), where \(q_\circ\ge0\), the factorial identities are
\[
\mathbb E(\mathsf N_\circ)_j=(m q_\circ)^j,
\]
\[
(\mathsf N_\circ)_j(\mathsf N_\circ)_k
=
\sum_{\ell=0}^{\min(j,k)}
\binom j\ell\binom k\ell\ell!\,
(\mathsf N_\circ)_{j+k-\ell},
\qquad j,k\in\mathbb N.
\]
The first follows by shifting the Poisson series. The second counts the overlap of two ordered selections of distinct indices. Substituting the first identity into the centered factorial statistic of \cref{obj:def:armwise-estimator}, for every real \(\zeta\), and applying the binomial theorem gives
\[
\mathbb E U_j(\mathsf N_\circ;\zeta)
=
\sum_{j_{\mathrm{raw}}=0}^{j}
\binom j{j_{\mathrm{raw}}}(-\zeta)^{j-j_{\mathrm{raw}}}
q_\circ^{j_{\mathrm{raw}}}
=(q_\circ-\zeta)^j.
\]
For the mixed second moment, substitution of the raw factorial product identity gives the finite triple sum
\[
\begin{aligned}
\mathbb E[U_j(\mathsf N_\circ;\zeta)U_k(\mathsf N_\circ;\zeta)]
={}&
\sum_{j_{\mathrm{raw}}=0}^{j}
\sum_{k_{\mathrm{raw}}=0}^{k}
\binom j{j_{\mathrm{raw}}}\binom k{k_{\mathrm{raw}}}
(-\zeta)^{j-j_{\mathrm{raw}}+k-k_{\mathrm{raw}}}\\
&\quad\times
\sum_{\ell=0}^{\min(j_{\mathrm{raw}},k_{\mathrm{raw}})}
\binom{j_{\mathrm{raw}}}\ell
\binom{k_{\mathrm{raw}}}\ell\ell!
\left(\frac{q_\circ}{m}\right)^\ell
q_\circ^{j_{\mathrm{raw}}+k_{\mathrm{raw}}-2\ell}.
\end{aligned}
\]
Fixing the overlap order and regrouping the finite sums factors this expression as
\[
\begin{aligned}
\mathbb E[U_j(\mathsf N_\circ;\zeta)U_k(\mathsf N_\circ;\zeta)]
={}&\sum_{\ell=0}^{\min(j,k)}
\ell!\left(\frac{q_\circ}{m}\right)^\ell\\
&\times\left\{
\sum_{j_{\mathrm{raw}}=\ell}^{j}
\binom j{j_{\mathrm{raw}}}\binom{j_{\mathrm{raw}}}\ell
(-\zeta)^{j-j_{\mathrm{raw}}}
q_\circ^{j_{\mathrm{raw}}-\ell}
\right\}\\
&\times\left\{
\sum_{k_{\mathrm{raw}}=\ell}^{k}
\binom k{k_{\mathrm{raw}}}\binom{k_{\mathrm{raw}}}\ell
(-\zeta)^{k-k_{\mathrm{raw}}}
q_\circ^{k_{\mathrm{raw}}-\ell}
\right\}.
\end{aligned}
\]
Terms with a raw order smaller than \(\ell\) vanish because the corresponding binomial coefficient is zero. For the remaining terms, the shifted binomial identity is
\[
\begin{aligned}
\sum_{j_{\mathrm{raw}}=\ell}^{j}
\binom j{j_{\mathrm{raw}}}\binom{j_{\mathrm{raw}}}\ell
(-\zeta)^{j-j_{\mathrm{raw}}}
q_\circ^{j_{\mathrm{raw}}-\ell}
&=
\binom j\ell
\sum_{i_{\mathrm{shift}}=0}^{j-\ell}
\binom{j-\ell}{i_{\mathrm{shift}}}
(-\zeta)^{j-\ell-i_{\mathrm{shift}}}
q_\circ^{i_{\mathrm{shift}}}\\
&=\binom j\ell(q_\circ-\zeta)^{j-\ell},
\qquad 0\le\ell\le j.
\end{aligned}
\]
The first equality substitutes \(j_{\mathrm{raw}}=\ell+i_{\mathrm{shift}}\) and uses
\[
\binom j{\ell+i_{\mathrm{shift}}}
\binom{\ell+i_{\mathrm{shift}}}\ell
=
\binom j\ell\binom{j-\ell}{i_{\mathrm{shift}}},
\qquad
 i_{\mathrm{shift}}\in\{0,\ldots,j-\ell\};
\]
the second is the binomial theorem. Applying this identity to both factors yields
\[
\mathbb E[U_j(\mathsf N_\circ;\zeta)U_k(\mathsf N_\circ;\zeta)]
=
\sum_{\ell=0}^{\min(j,k)}
\binom j\ell\binom k\ell\ell!
\left(\frac{q_\circ}{m}\right)^\ell
(q_\circ-\zeta)^{j+k-2\ell}.
\]
Taking \(k=j\) therefore gives
\[
\mathbb E U_j(\mathsf N_\circ;\zeta)^2
=
\sum_{\ell=0}^{j}
\binom j\ell^2\ell!
\left(\frac{q_\circ}{m}\right)^\ell
(q_\circ-\zeta)^{2j-2\ell}.
\]
Thus, conditional on a pilot in \(\mathcal G_x\), for \(0\le j\le D\),
\[
\mathbb E_{\mathrm{aux}}
\left[
\left.
\frac{U_j(N_{ay,x};b_{ay,x})^2}{r_{ay,x}^{2j}}
\right|N'
\right]
\le
\sum_{\ell=0}^{j}\frac{j^{2\ell}}{\ell!}
\left(\frac8L\right)^\ell
\le e^{8j^2/L}.
\]
Here \(N'\) denotes the complete pilot-count array. Evaluation-coordinate independence therefore gives
\[
\mathbb E_{\mathrm{aux}}[Z_x\mid N']
=P_{x,K_n}(q_x),
\]
\[
\mathbb E_{\mathrm{aux}}
[(Z_x-F_\epsilon(b_x))^2\mid N']
\le
\left(\sum_{\boldsymbol\alpha}|\beta_{x,\boldsymbol\alpha}|\right)^2
e^{32D^2/L}
\quad\text{on }\mathcal G_x.
\]
For the last inequality, each normalized factorial product has second moment at most \(e^{32D^2/L}\); Cauchy--Schwarz bounds each cross term in the polynomial square by this same envelope.

Choose the numerical constants
\[
\kappa=\frac1{17952},
\qquad
C_{\mathrm{fac}}=\exp\left(1632+\frac1{16}\right).
\]
This choice satisfies
\[
304\kappa+256\kappa^2\le\frac1{16}.
\]
The degree definition gives, for every \(d\ge1\),
\[
\frac{\kappa L}{2}\le K_n\le2+\kappa L.
\]
For the lower bound, split according as \(\kappa L\le4\) or \(\kappa L>4\), and use respectively \(K_n\ge2\) or
\(\lfloor\kappa L\rfloor>\kappa L-1\). Also,
\[
K_n^2\le8+2\kappa^2L^2.
\]
Combining these inequalities with the coefficient estimate gives
\[
\begin{aligned}
\left(\sum_{\boldsymbol\alpha}|\beta_{x,\boldsymbol\alpha}|\right)^2
e^{32D^2/L}
&\le
(S_x^{\mathrm{aw}})^2
e^{304K_n+128K_n^2/L}\\
&\le
(S_x^{\mathrm{aw}})^2e^{1632+L/16}\\
&=
C_{\mathrm{fac}}d^{1/16}(S_x^{\mathrm{aw}})^2.
\end{aligned}
\]
Thus the degree tuning is uniform over the entire auxiliary alphabet range.
% lean: CausalSmith.Stat.OptvalueVanishingoverlapRate.pilotFactorialCellValue_square_envelope
% lean: CausalSmith.Stat.OptvalueVanishingoverlapRate.pilotJacksonCoeff_square_tuning_exists

\textit{4. Clipping on successful pilots.}
On \(\mathcal G_x\), the polynomial oscillation bound places \(P_{x,K_n}(q_x)\) in the clipping interval, since \(d^{1/4}\ge1\). Projection onto that interval consequently gives
\[
\mathbb E_{\mathrm{aux}}
[(T_x-P_{x,K_n}(q_x))^2\mid N']
\le
C_{\mathrm{fac}}d^{1/16}(S_x^{\mathrm{aw}})^2.
\]
Indeed, projection contracts distance from \(P_{x,K_n}(q_x)\), and
\[
\mathbb E_{\mathrm{aux}}
[(Z_x-P_{x,K_n}(q_x))^2\mid N']
\le
\mathbb E_{\mathrm{aux}}
[(Z_x-F_\epsilon(b_x))^2\mid N'].
\]
For the clipping bias, the scalar projection inequality
\[
|\Pi_{[-\mathfrak u_{\mathrm{clip}},\mathfrak u_{\mathrm{clip}}]}(\mathfrak v_{\mathrm{clip}})-\mathfrak v_{\mathrm{clip}}|
\le\frac{\mathfrak v_{\mathrm{clip}}^2}{\mathfrak u_{\mathrm{clip}}},
\qquad \mathfrak u_{\mathrm{clip}}>0,\quad
\mathfrak v_{\mathrm{clip}}\in\mathbb R,
\]
follows by considering \(|\mathfrak v_{\mathrm{clip}}|\le\mathfrak u_{\mathrm{clip}}\) and \(|\mathfrak v_{\mathrm{clip}}|>\mathfrak u_{\mathrm{clip}}\). Applying it with clipping radius \(d^{1/4}S_x^{\mathrm{aw}}>0\), and using factorial unbiasedness, yields
\[
\left|
\mathbb E_{\mathrm{aux}}[T_x\mid N']-P_{x,K_n}(q_x)
\right|
\le C_{\mathrm{fac}}d^{-3/16}S_x^{\mathrm{aw}}
\quad\text{on }\mathcal G_x.
\]
% lean: CausalSmith.Stat.OptvalueVanishingoverlapRate.pilotClippedCellValue_tuning_exists

\textit{5. Approximation error in occupied cells.}
For \(p_x>0\), define
\[
\rho_{1x}=\pi_x,\qquad \rho_{0x}=1-\pi_x.
\]
By \cref{obj:def:observed-class,obj:ass:shrinking-overlap},
\[
s(q_x)=p_x,\qquad
s_a(q_x)=p_x\rho_{ax}>0,\qquad
D_a(q_x)=s_a(q_x),\qquad
\rho_{ax}\ge\epsilon.
\]
It follows from \cref{obj:def:armwise-extension,obj:def:observed-value} that
\[
F_\epsilon(q_x)=p_x\max_a\mu_{ax},
\qquad
\Psi(\mathbb P)=\sum_{x\in[d]}F_\epsilon(q_x).
\]
The same cell identity holds when \(p_x=0\), because then \(q_x=0\).

We record the kernel moments used in the approximation. Define, for \(t\in[-\pi,\pi]\), the continuously extended sine ratio and its fourth-power mass by
\[
\mathcal S_{K_n}(t)=
\begin{cases}
\displaystyle\frac{\sin(K_nt/2)}{\sin(t/2)},&t\ne0,\\[4pt]
K_n,&t=0,
\end{cases}
\qquad
\mathcal A_{K_n}=\int_{-\pi}^{\pi}\mathcal S_{K_n}(t)^4\,dt.
\]
The finite cosine expansion is
\[
\mathcal S_{K_n}(t)^2
=K_n+2\sum_{j=1}^{K_n-1}(K_n-j)\cos(jt).
\]
Indeed, multiplying the right side by \(2(1-\cos t)\) and using the cosine product identity telescopes to \(2(1-\cos(K_nt))\); division by \(4\sin^2(t/2)\) proves the identity for \(t\ne0\), and continuity supplies its value at zero. Cosine orthogonality and the finite square sum give
\[
\sum_{j=1}^{K_n-1}\{2(K_n-j)\}^2
=4\sum_{j=1}^{K_n-1}j^2
=\frac23(K_n-1)K_n(2K_n-1),
\]
where the square-sum formula follows by summing the successive differences of cubes. Consequently,
\[
\begin{aligned}
\mathcal A_{K_n}
&=2\pi K_n^2+
\pi\sum_{j=1}^{K_n-1}\{2(K_n-j)\}^2\\
&=\frac{2\pi}{3}K_n(2K_n^2+1)
\ge\frac{4\pi}{3}K_n^3,
\qquad
J_{K_n}(t)=\frac{\mathcal S_{K_n}(t)^4}{\mathcal A_{K_n}}.
\end{aligned}
\]
The same cosine expansion, integrated without squaring, yields
\[
\int_{-\pi}^{\pi}\mathcal S_{K_n}(t)^2\,dt=2\pi K_n.
\]
For \(|t|\le\pi\), the sine lower bound gives
\[
|t|\le\pi|\sin(t/2)|,
\qquad
t^2\le\pi^2\sin^2(t/2).
\]
For \(t\ne0\), it follows that
\[
\begin{aligned}
t^2\mathcal S_{K_n}(t)^4
&\le\pi^2\sin^2(K_nt/2)\mathcal S_{K_n}(t)^2\\
&\le\pi^2\mathcal S_{K_n}(t)^2.
\end{aligned}
\]
At \(t=0\), this inequality follows directly from its zero left side. Integrating and normalizing therefore gives
\[
\int_{-\pi}^{\pi}t^2\mathcal S_{K_n}(t)^4\,dt
\le2\pi^3K_n,
\qquad
\int_{-\pi}^{\pi}t^2J_{K_n}(t)\,dt
\le\frac{24}{K_n^2}\le\frac{64}{K_n^2},
\]
using the mass lower bound and \(\pi^2\le16\).
Finally, the nonnegative square \((K_n|t|-2)^2\) implies
\[
|t|\le\frac{K_n}{4}t^2+\frac1{K_n}.
\]
Since the kernel is nonnegative and integrates to one, we obtain
\[
\begin{aligned}
\int_{-\pi}^{\pi}|t|J_{K_n}(t)\,dt
&\le\frac{K_n}{4}
\int_{-\pi}^{\pi}t^2J_{K_n}(t)\,dt+\frac1{K_n}\\
&\le\frac7{K_n}\le\frac{32}{K_n}.
\end{aligned}
\]
% lean: Causalean.Mathlib.Analysis.JacksonApproximation.jrawMass_eq
% lean: Causalean.Mathlib.Analysis.JacksonApproximation.jackson_second_moment
% lean: Causalean.Mathlib.Analysis.JacksonApproximation.jackson_first_moment
For \(q_{ay,x}=b_{ay,x}+r_{ay,x}\cos\theta_{ay,x}\), define
\[
\theta_{ay,x}
=
\arccos\left(\frac{q_{ay,x}-b_{ay,x}}{r_{ay,x}}\right)
\quad\text{on }\mathcal G_x.
\]
Taylor's formula gives
\[
\begin{aligned}
r_{ay,x}
|\cos(\theta_{ay,x}+t)-\cos\theta_{ay,x}|
\le{}&
\sqrt{(q_{ay,x}-\ell_{ay,x})(u_{ay,x}-q_{ay,x})}\,|t|\\
&+\frac{r_{ay,x}}2t^2.
\end{aligned}
\]
Applying \cref{obj:lem:armwise-modulus} at \(q_x\) inside the tensor convolution therefore gives
\[
\begin{aligned}
|P_{x,K_n}(q_x)-F_\epsilon(q_x)|
\le64\sum_{a=0}^1
\left(1+\frac{p_x}{s_a(q_x)}\right)
\left\{
\frac{\sum_y\sqrt{(q_{ay,x}-\ell_{ay,x})(u_{ay,x}-q_{ay,x})}}{K_n}
+\frac{R_{a,x}}{K_n^2}
\right\}.
\end{aligned}
\]

The localized coordinate formulas imply
\[
h_{ay,x}
\le H_0(H_0+2)
\left(\sqrt{q_{ay,x}\delta_{\mathrm p}}+\delta_{\mathrm p}\right).
\]
To see this, the lower-endpoint inequality implies
\[
c_{ay,x}-H_0\left(\sqrt{c_{ay,x}\delta_{\mathrm p}}+\delta_{\mathrm p}\right)
\le q_{ay,x}.
\]
Solving the resulting quadratic inequality in \(\sqrt{c_{ay,x}}\) gives
\[
\sqrt{c_{ay,x}\delta_{\mathrm p}}
\le
\sqrt{q_{ay,x}\delta_{\mathrm p}}+(H_0+1)\delta_{\mathrm p},
\]
which yields the stated half-width bound.

There is also the boundary-sensitive estimate
\[
\sqrt{(q_{ay,x}-\ell_{ay,x})(u_{ay,x}-q_{ay,x})}
\le2H_0(H_0+2)\sqrt{q_{ay,x}\delta_{\mathrm p}}.
\]
For \(q_{ay,x}\le\delta_{\mathrm p}\), bound the left distance by \(q_{ay,x}\), the right distance by \(2h_{ay,x}\), and the half width by
\(2H_0(H_0+2)\delta_{\mathrm p}\). For \(q_{ay,x}>\delta_{\mathrm p}\), the product of the two distances is at most \(h_{ay,x}^2\), and the half-width bound applies. Summing coordinate radii also gives
\[
R_{a,x}
\le2H_0(H_0+2)
\left(\sqrt{s_a(q_x)\delta_{\mathrm p}}+\delta_{\mathrm p}\right).
\]

For \(d\ge2\), the mass bound in \cref{obj:lem:armwise-modulus} supplies
\[
\sum_{a=0}^1\frac{p_x}{s_a(q_x)}
\sum_{y=0}^1\sqrt{q_{ay,x}}
\le2\sqrt2\sqrt{\frac{p_x}{\epsilon}}.
\]
For \(d=1\), the unique cell has \(p_x=1\), and the same bound follows directly from the two arm masses. For each arm, Cauchy--Schwarz and \(s_a(q_x)\ge\epsilon p_x>0\) give
\[
\sum_{y=0}^1\sqrt{q_{ay,x}}\le\sqrt{2s_a(q_x)},
\qquad
\frac{p_x}{s_a(q_x)}\sum_{y=0}^1\sqrt{q_{ay,x}}
\le\frac{\sqrt2\,p_x}{\sqrt{s_a(q_x)}}
\le\sqrt{\frac{2p_x}{\epsilon}}.
\]
Adding the two arm inequalities proves the displayed mass bound also in this case.
% lean: CausalSmith.Stat.OptvalueVanishingoverlapRate.observed_armwise_sqrt_bound
Together with \(1+p_x/s_a(q_x)\le2p_x/s_a(q_x)\), this bound now yields
\[
\begin{aligned}
|P_{x,K_n}(q_x)-F_\epsilon(q_x)|
\le512H_0(H_0+2)
\left\{
\frac{\sqrt2}{K_n}\sqrt{\frac{p_xL}{m\epsilon}}
+
\frac1{K_n^2}
\left(\sqrt{\frac{p_xL}{m\epsilon}}+\frac L{m\epsilon}\right)
\right\}.
\end{aligned}
\]
Define
\[
\eta_x=
\sqrt{\frac{p_x}{m\epsilon L}}+\frac1{m\epsilon L},
\]
\[
C_{\mathrm{occ}}
=
512H_0(H_0+2)\sqrt2
\left(\frac1{\kappa/2}+\frac1{(\kappa/2)^2}\right).
\]
Using \(K_n\ge(\kappa/2)L\) and \(L\ge1\), we conclude that
\[
|P_{x,K_n}(q_x)-F_\epsilon(q_x)|
\le C_{\mathrm{occ}}\eta_x
\quad\text{on }\mathcal G_x,\quad p_x>0.
\]
Combining this estimate with the clipping bounds and integrating over the pilot gives
\[
\mathbb E_{\mathrm{aux}}
[(T_x-F_\epsilon(q_x))^2\mathbf1_{\mathcal G_x}]
\le
2C_{\mathrm{fac}}d^{1/16}
\mathbb E_{\mathrm{aux}}(S_x^{\mathrm{aw}})^2
+2C_{\mathrm{occ}}^2\eta_x^2,
\]
\[
\left|
\mathbb E_{\mathrm{aux}}
[(T_x-F_\epsilon(q_x))\mathbf1_{\mathcal G_x}]
\right|
\le
C_{\mathrm{occ}}\eta_x
+
C_{\mathrm{fac}}d^{-3/16}
\mathbb E_{\mathrm{aux}}S_x^{\mathrm{aw}}.
\]
The squared bound uses
\((v+w)^2\le2v^2+2w^2\); the mean bound uses the triangle inequality.
% lean: CausalSmith.Stat.OptvalueVanishingoverlapRate.observed_pilotJackson_mean_bias_normalized_le
% lean: CausalSmith.Stat.OptvalueVanishingoverlapRate.markedPoissonCellValue_occupied_good_moments_tuning_exists

\textit{6. Null cells and uniform successful-pilot bounds.}
If \(p_x=0\), every pilot and evaluation count in that cell is zero almost surely. Hence
\[
q_x=0,\qquad
b_{ay,x}=r_{ay,x}=\frac{H_0L}{2m},\qquad
\ell_{ay,x}=0,
\qquad
\mathbb P_{\mathrm{aux}}(\mathcal G_x)=1.
\]
At the lower corner, the factorial statistic equals the polynomial value:
\[
Z_x=P_{x,K_n}(0).
\]
This follows either by substituting zero counts in the factorial expansion or by its expectation identity, since the evaluation counts are constant. The polynomial oscillation bound places this value inside the clipping interval, so
\[
T_x=P_{x,K_n}(0).
\]
From \cref{obj:def:armwise-extension},
\[
0\le F_\epsilon(v)\le s(v),\qquad v\in\mathbb R_+^4.
\]
At the null corner, a displaced coordinate in the convolution is
\(r_{ay,x}(1-\cos t)\le r_{ay,x}t^2/2\). The second kernel moment thus gives
\[
|T_x-F_\epsilon(0)|
\le\frac{32}{K_n^2}\sum_{a,y}r_{ay,x}
=\frac{64H_0L}{mK_n^2}.
\]
Define
\[
C_{\mathrm{null}}=\frac{64H_0}{(\kappa/2)^2},
\qquad
C_{\mathrm{app}}=C_{\mathrm{occ}}+C_{\mathrm{null}}.
\]
Since \(K_n\ge(\kappa/2)L\) and \(\epsilon\le1\), the null-cell error is at most \(C_{\mathrm{null}}\eta_x\).

Therefore, for every cell, occupied or null, the successful-pilot estimates take the common form
\[
\mathbb E_{\mathrm{aux}}
[(T_x-F_\epsilon(q_x))^2\mathbf1_{\mathcal G_x}]
\le
C_{\mathrm v}d^{1/16}
\mathbb E_{\mathrm{aux}}(S_x^{\mathrm{aw}})^2
+C_{\mathrm e}\eta_x^2,
\]
\[
\left|
\mathbb E_{\mathrm{aux}}
[(T_x-F_\epsilon(q_x))\mathbf1_{\mathcal G_x}]
\right|
\le
C_{\mathrm b}
\left(\eta_x+d^{-3/16}\mathbb E_{\mathrm{aux}}S_x^{\mathrm{aw}}\right),
\]
where the universal constants are
\[
C_{\mathrm v}=2C_{\mathrm{fac}},\qquad
C_{\mathrm e}=2C_{\mathrm{app}}^2,\qquad
C_{\mathrm b}=C_{\mathrm{app}}+C_{\mathrm{fac}}.
\]
These enlargements use the nonnegativity of the clipping scales and of \(\eta_x\).
% lean: CausalSmith.Stat.OptvalueVanishingoverlapRate.markedPoissonCellValue_null_good_moments_le
% lean: CausalSmith.Stat.OptvalueVanishingoverlapRate.markedPoissonStatistic_active_rate_tuning_exists

\textit{7. Unconditional clipping-scale moments.}
For each pilot, the coordinate inequalities \(r_{ay,x}\le h_{ay,x}\) and \(c_{ay,x}\le b_{ay,x}\) imply
\[
R_{a,x}
\le2H_0\left(\sqrt{s_a(b_x)\delta_{\mathrm p}}+\delta_{\mathrm p}\right).
\]
Using \(D_a(b_x)\ge s_a(b_x)\) and
\(D_a(b_x)\ge\epsilon s(b_x)\), separately according as
\(s_a(b_x)\ge\epsilon s(b_x)\) or
\(s_a(b_x)<\epsilon s(b_x)\), gives
\[
\frac{s(b_x)}{D_a(b_x)}R_{a,x}
\le
2H_0\left(
\sqrt{\frac{s(b_x)\delta_{\mathrm p}}{\epsilon}}
+\frac{\delta_{\mathrm p}}{\epsilon}
\right).
\]
The unweighted radius satisfies this same bound, since
\(s_a(b_x)\le s(b_x)\) and \(\epsilon\le1\). Consequently,
\[
S_x^{\mathrm{aw}}
\le16H_0\left(
\sqrt{\frac{s(b_x)\delta_{\mathrm p}}{\epsilon}}
+\frac{\delta_{\mathrm p}}{\epsilon}
\right),
\]
\[
(S_x^{\mathrm{aw}})^2
\le512H_0^2\left(
\frac{s(b_x)\delta_{\mathrm p}}{\epsilon}
+\frac{\delta_{\mathrm p}^2}{\epsilon^2}
\right).
\]

The midpoint formulas also give
\[
b_{ay,x}\le(1+H_0)c_{ay,x}+2H_0\delta_{\mathrm p},
\qquad
s(b_x)\le(1+H_0)\sum_{a,y}c_{ay,x}+8H_0\delta_{\mathrm p}.
\]
For example, use \(b_{ay,x}\le c_{ay,x}+h_{ay,x}\) and
\(\sqrt{c_{ay,x}\delta_{\mathrm p}}\le c_{ay,x}+\delta_{\mathrm p}\).
Because \(\mathbb E_{\mathrm{aux}}c_{ay,x}=q_{ay,x}\), integration yields
\[
\mathbb E_{\mathrm{aux}}(S_x^{\mathrm{aw}})^2
\le512H_0^2
\left\{
(1+H_0)\frac{p_x\delta_{\mathrm p}}{\epsilon}
+8H_0\frac{\delta_{\mathrm p}^2}{\epsilon}
+\frac{\delta_{\mathrm p}^2}{\epsilon^2}
\right\}.
\]
Define, for every cell,
\[
\mathfrak a_x
=
\frac{p_xL}{m\epsilon}
+\frac{L^2}{m^2\epsilon^2},
\qquad
C_{\mathrm s}=512H_0^2(1+9H_0).
\]
Since \(\epsilon\le1\),
\[
\mathbb E_{\mathrm{aux}}(S_x^{\mathrm{aw}})^2
\le C_{\mathrm s}\mathfrak a_x.
\]
% lean: CausalSmith.Stat.OptvalueVanishingoverlapRate.markedPoissonClippingScale_integral_sq_le

\textit{8. Failed-pilot moments.}
Define the nonnegative coordinate scores by
\[
e_{ay,x}^{\mathrm p}
=
|c_{ay,x}-q_{ay,x}|+h_{ay,x},
\qquad
e_x^{\mathrm p}=\sum_{a,y}e_{ay,x}^{\mathrm p}.
\]
We need both a probability bound and moments retaining the cell mass:
\[
\mathbb P_{\mathrm{aux}}(\mathcal G_x^c)\le8e^{-40L},
\]
\[
\mathbb E_{\mathrm{aux}}(e_{ay,x}^{\mathrm p})^4
\le
\left(\frac{H_0}{1024}\right)^4
10^{20}
\left(\sqrt{q_{ay,x}\delta_{\mathrm p}}+\delta_{\mathrm p}\right)^4,
\]
\[
\mathbb E_{\mathrm{aux}}
[e_x^{\mathrm p}\mathbf1_{\mathcal G_x^c}]
\le
\frac{H_0}{1024}\,10^{160}e^{-20L}
\left(\sqrt{p_x\delta_{\mathrm p}}+\delta_{\mathrm p}\right).
\]

Here are the underlying calculations. For any
\(\lambda_\circ\ge0\) and
\(\mathsf N_\circ\sim\operatorname{Pois}(\lambda_\circ)\), define
\[
\mathfrak s_\circ=\sqrt{\lambda_\circ L}+L,\qquad
\mathfrak e_\circ
=
|\mathsf N_\circ-\lambda_\circ|
+1024\left(\sqrt{\mathsf N_\circ L}+L\right),
\qquad
\vartheta_\circ=\frac1{2\mathfrak s_\circ}.
\]
The centered Poisson generating function is
\[
\mathbb E e^{t(\mathsf N_\circ-\lambda_\circ)}
=
\exp\{\lambda_\circ(e^t-1-t)\}.
\]
If \(\lambda_\circ=0\), then \(\mathsf N_\circ=0\) almost surely, and both tail events below have probability zero. For \(\lambda_\circ>0\), define the deviation threshold and exponential tilt at an arbitrary nonnegative level by
\[
z_{\mathrm{tail}}\in[0,\infty),\qquad
\mathfrak x_{\mathrm{up}}
=\sqrt{2\lambda_\circ z_{\mathrm{tail}}}+2z_{\mathrm{tail}},
\qquad
\vartheta_{\mathrm{up}}
=\log\left(1+\frac{\mathfrak x_{\mathrm{up}}}{\lambda_\circ}\right)\ge0.
\]
Exponential Markov inequality and the generating function give
\[
\begin{aligned}
\Pr\{\mathsf N_\circ-\lambda_\circ>\mathfrak x_{\mathrm{up}}\}
&\le
\exp\left\{
-\vartheta_{\mathrm{up}}\mathfrak x_{\mathrm{up}}
+\lambda_\circ(e^{\vartheta_{\mathrm{up}}}-1-\vartheta_{\mathrm{up}})
\right\}\\
&=
\exp\left\{
-\bigl[(\lambda_\circ+\mathfrak x_{\mathrm{up}})
\vartheta_{\mathrm{up}}-\mathfrak x_{\mathrm{up}}\bigr]
\right\}.
\end{aligned}
\]
The logarithmic lower bound follows from
\[
\frac{d}{d\xi_{\mathrm{log}}}
\left\{\log(1+\xi_{\mathrm{log}})
-\frac{2\xi_{\mathrm{log}}}{\xi_{\mathrm{log}}+2}\right\}
=
\frac{\xi_{\mathrm{log}}^2}
{(1+\xi_{\mathrm{log}})(\xi_{\mathrm{log}}+2)^2}
\ge0,
\qquad \xi_{\mathrm{log}}\in[0,\infty),
\]
since the difference is zero at \(\xi_{\mathrm{log}}=0\). Applying it at \(\mathfrak x_{\mathrm{up}}/\lambda_\circ\) gives
\[
\vartheta_{\mathrm{up}}
\ge\frac{2\mathfrak x_{\mathrm{up}}}{2\lambda_\circ+\mathfrak x_{\mathrm{up}}},
\qquad
(\lambda_\circ+\mathfrak x_{\mathrm{up}})\vartheta_{\mathrm{up}}
-\mathfrak x_{\mathrm{up}}
\ge
\frac{\mathfrak x_{\mathrm{up}}^2}
{2\lambda_\circ+\mathfrak x_{\mathrm{up}}}
\ge z_{\mathrm{tail}}.
\]
The last inequality follows by expanding the threshold:
\[
\mathfrak x_{\mathrm{up}}^2
-z_{\mathrm{tail}}(2\lambda_\circ+\mathfrak x_{\mathrm{up}})
=
3z_{\mathrm{tail}}\sqrt{2\lambda_\circ z_{\mathrm{tail}}}
+2z_{\mathrm{tail}}^2\ge0.
\]
Thus the upper-tail probability is at most \(e^{-z_{\mathrm{tail}}}\). Taking \(z_{\mathrm{tail}}=40L\), and applying the lower-tail exponential Markov inequality with
\(e^{-t}-1+t\le t^2/2\), \(t\ge0\), gives
\[
\Pr\!\left\{
\mathsf N_\circ-\lambda_\circ>
\sqrt{80\lambda_\circ L}+80L
\right\}\le e^{-40L},
\qquad
\Pr\!\left\{
\lambda_\circ-\mathsf N_\circ>
\sqrt{80\lambda_\circ L}
\right\}\le e^{-40L}.
\]
The event
\[
\left\{
|\mathsf N_\circ-\lambda_\circ|
>\frac{1024}{4}\left(\sqrt{\mathsf N_\circ L}+L\right)
\right\}
\]
is contained in the union of these two tail events: on their complement, substitute
\(\lambda_\circ\le\mathsf N_\circ+|\mathsf N_\circ-\lambda_\circ|\) and square the nonnegative terms. Since \(H_0\ge1024\), failure of a pilot coordinate implies the corresponding event above. A union bound over the four coordinates proves the probability estimate.

For the moments, square-root subadditivity and
\(\sqrt{|\mathsf N_\circ-\lambda_\circ|L}
\le\{|\mathsf N_\circ-\lambda_\circ|+L\}/2\) imply
\[
\mathfrak e_\circ
\le1536\left(\mathfrak s_\circ+
|\mathsf N_\circ-\lambda_\circ|\right).
\]
Also \(\mathfrak s_\circ\ge1\),
\(\lambda_\circ\le\mathfrak s_\circ^2\), and
\(\lambda_\circ\vartheta_\circ^2\le1/4\). Using
\(|e^t-1-t|\le t^2\) for \(|t|\le1\) in the generating function therefore gives
\[
\mathbb E\exp\left(\frac{\mathfrak e_\circ}{4096\mathfrak s_\circ}\right)
\le
e^{1/2}\left(
\mathbb E e^{\vartheta_\circ(\mathsf N_\circ-\lambda_\circ)}
+
\mathbb E e^{-\vartheta_\circ(\mathsf N_\circ-\lambda_\circ)}
\right)
\le2e^{3/4}\le8.
\]
Thus, for \(j=2,4\), the inequality \(v^j\le j!e^v\), \(v\ge0\), yields
\[
\mathbb E\mathfrak e_\circ^j
\le8j!\,4096^j\mathfrak s_\circ^j
\le10^{4j+4}\mathfrak s_\circ^j.
\]
Substituting \(\lambda_\circ=mq_{ay,x}\) gives the displayed fourth-moment bound, because
\[
e_{ay,x}^{\mathrm p}
\le\frac{H_0}{1024m}\mathfrak e_\circ.
\]
We next localize the scalar moments before summing over failing coordinates. For the scalar Poisson count above, define
\[
\Delta_\circ=|\mathsf N_\circ-\lambda_\circ|,
\qquad
\mathcal B_\circ=
\left\{
\Delta_\circ>256\left(\sqrt{\mathsf N_\circ L}+L\right)
\right\}.
\]
Applying the exponential Markov calculations at level \(80L\) gives
\[
\Pr\!\left\{
\mathsf N_\circ-\lambda_\circ>
\sqrt{160\lambda_\circ L}+160L
\right\}\le e^{-80L},
\qquad
\Pr\!\left\{
\lambda_\circ-\mathsf N_\circ>
\sqrt{160\lambda_\circ L}
\right\}\le e^{-80L}.
\]
To compare \(\mathcal B_\circ\) with these events, first suppose \(\mathsf N_\circ\ge\lambda_\circ\). On \(\mathcal B_\circ\),
\[
\sqrt{160\lambda_\circ L}+160L
\le16\sqrt{\mathsf N_\circ L}+160L
\le256\left(\sqrt{\mathsf N_\circ L}+L\right)
<\Delta_\circ.
\]
If \(\mathsf N_\circ<\lambda_\circ\), then \(\lambda_\circ=\mathsf N_\circ+\Delta_\circ\). The defining inequality of \(\mathcal B_\circ\), and monotonicity of the quadratic above that threshold, give
\[
\begin{aligned}
\Delta_\circ^2-160L\Delta_\circ-160\mathsf N_\circ L
>{}&65376\mathsf N_\circ L
+90112L\sqrt{\mathsf N_\circ L}+24576L^2>0.
\end{aligned}
\]
Thus \(\Delta_\circ>\sqrt{160\lambda_\circ L}\). The two cases prove
\[
\Pr(\mathcal B_\circ)\le2e^{-80L}.
\]
The exponential score bound above also yields, for every integer \(1\le j_{\mathrm{mom}}\le8\),
\[
\mathbb E\mathfrak e_\circ^{j_{\mathrm{mom}}}
\le8j_{\mathrm{mom}}!\,4096^{j_{\mathrm{mom}}}
\mathfrak s_\circ^{j_{\mathrm{mom}}}
\le10^{4j_{\mathrm{mom}}+4}\mathfrak s_\circ^{j_{\mathrm{mom}}}.
\]
For \(j_{\mathrm{mom}}\in\{1,2,4\}\), define the positive balancing coefficient by
\[
\upsilon_{\circ,j_{\mathrm{mom}}}
=\frac{e^{-40L}}{\mathfrak s_\circ^{j_{\mathrm{mom}}}}.
\]
The nonnegative square
\((\upsilon_{\circ,j_{\mathrm{mom}}}\mathfrak e_\circ^{j_{\mathrm{mom}}}-1)^2\), on \(\mathcal B_\circ\), proves the pointwise inequality
\[
\mathfrak e_\circ^{j_{\mathrm{mom}}}\mathbf1_{\mathcal B_\circ}
\le\frac12\left(
\upsilon_{\circ,j_{\mathrm{mom}}}\mathfrak e_\circ^{2j_{\mathrm{mom}}}
+\upsilon_{\circ,j_{\mathrm{mom}}}^{-1}\mathbf1_{\mathcal B_\circ}
\right).
\]
Outside \(\mathcal B_\circ\), its left side is zero and its right side is nonnegative. Integration therefore gives the localized scalar bound
\[
\begin{aligned}
\mathbb E\!\left[
\mathfrak e_\circ^{j_{\mathrm{mom}}}\mathbf1_{\mathcal B_\circ}
\right]
&\le\frac{10^{8j_{\mathrm{mom}}+4}+2}{2}
 e^{-40L}\mathfrak s_\circ^{j_{\mathrm{mom}}}\\
&\le10^{8j_{\mathrm{mom}}+8}
 e^{-40L}\mathfrak s_\circ^{j_{\mathrm{mom}}}.
\end{aligned}
\]
% lean: Causalean.Stat.Concentration.PoissonSelfNormalized.poisson_weighted_bad_moment

For each \((a,y)\in\{0,1\}^2\) and each cell \(x\), define the raw score, raw local scale, and raw failure event by
\[
\mathfrak e_{ay,x}
=|N'_{ay,x}-mq_{ay,x}|
+1024\left(\sqrt{N'_{ay,x}L}+L\right),
\qquad
\mathfrak s_{ay,x}=\sqrt{mq_{ay,x}L}+L,
\]
\[
\mathcal B_{ay,x}
=\left\{
|N'_{ay,x}-mq_{ay,x}|>
256\left(\sqrt{N'_{ay,x}L}+L\right)
\right\},
\qquad
\mathcal B_x^{\cup}=\bigcup_{a,y}\mathcal B_{ay,x}.
\]
The interval-failure criterion and \(H_0\ge1024\) imply
\[
\mathcal G_x^c\subseteq\mathcal B_x^{\cup},
\qquad
e_{ay,x}^{\mathrm p}
\le\frac{H_0}{1024m}\mathfrak e_{ay,x}.
\]
For \(j_{\mathrm{mom}}\in\{1,2,4\}\), consider the contribution of one coordinate's score on another coordinate's failure event. If \((a,y)=(a',y')\), the localized scalar bound gives
\[
\mathbb E_{\mathrm{aux}}\!\left[
\mathfrak e_{ay,x}^{j_{\mathrm{mom}}}\mathbf1_{\mathcal B_{ay,x}}
\right]
\le10^{8j_{\mathrm{mom}}+8}e^{-40L}
\mathfrak s_{ay,x}^{j_{\mathrm{mom}}}.
\]
If \((a,y)\ne(a',y')\), independence of the pilot coordinates factors the expectation:
\[
\begin{aligned}
\mathbb E_{\mathrm{aux}}\!\left[
\mathfrak e_{ay,x}^{j_{\mathrm{mom}}}\mathbf1_{\mathcal B_{a'y',x}}
\right]
&=\mathbb E_{\mathrm{aux}}\mathfrak e_{ay,x}^{j_{\mathrm{mom}}}
\Pr(\mathcal B_{a'y',x})\\
&\le2\cdot10^{4j_{\mathrm{mom}}+4}e^{-40L}
\mathfrak s_{ay,x}^{j_{\mathrm{mom}}},
\end{aligned}
\]
where the scalar probability bound was enlarged from \(2e^{-80L}\) to \(2e^{-40L}\). Hence both cases satisfy
\[
\mathbb E_{\mathrm{aux}}\!\left[
\mathfrak e_{ay,x}^{j_{\mathrm{mom}}}\mathbf1_{\mathcal B_{a'y',x}}
\right]
\le
\left(10^{8j_{\mathrm{mom}}+8}+2\cdot10^{4j_{\mathrm{mom}}+4}\right)
e^{-40L}\mathfrak s_{ay,x}^{j_{\mathrm{mom}}}.
\]
The union indicator and the finite-sum power inequality give the pointwise bound
\[
\left(\sum_{a,y}\mathfrak e_{ay,x}\right)^{j_{\mathrm{mom}}}
\mathbf1_{\mathcal B_x^{\cup}}
\le4^{j_{\mathrm{mom}}-1}
\sum_{a',y'}\sum_{a,y}
\mathfrak e_{ay,x}^{j_{\mathrm{mom}}}
\mathbf1_{\mathcal B_{a'y',x}}.
\]
To sum the local scales, Cauchy--Schwarz and \(\sum_{a,y}q_{ay,x}=p_x\) give
\[
\sum_{a,y}\mathfrak s_{ay,x}
\le2\sqrt{mp_xL}+4L
\le4\left(\sqrt{mp_xL}+L\right),
\qquad
\sum_{a,y}\mathfrak s_{ay,x}^{j_{\mathrm{mom}}}
\le4\left(\sum_{a,y}\mathfrak s_{ay,x}\right)^{j_{\mathrm{mom}}}.
\]
Integrating the pointwise bound and inserting the two coordinate cases therefore yields
\[
\begin{aligned}
\mathbb E_{\mathrm{aux}}\!\left[
\left(\sum_{a,y}\mathfrak e_{ay,x}\right)^{j_{\mathrm{mom}}}
\mathbf1_{\mathcal B_x^{\cup}}
\right]
\le{}&4^{2j_{\mathrm{mom}}+1}
\left(10^{8j_{\mathrm{mom}}+8}+2\cdot10^{4j_{\mathrm{mom}}+4}\right)
e^{-40L}\left(\sqrt{mp_xL}+L\right)^{j_{\mathrm{mom}}}\\
\le{}&10^{80(j_{\mathrm{mom}}+1)}e^{-20L}
\left(\sqrt{mp_xL}+L\right)^{j_{\mathrm{mom}}}.
\end{aligned}
\]
The last enlargement uses, for the three moment orders,
\[
4^{2j_{\mathrm{mom}}+1}
\left(10^{8j_{\mathrm{mom}}+8}+2\cdot10^{4j_{\mathrm{mom}}+4}\right)
\le10^{80(j_{\mathrm{mom}}+1)},
\qquad e^{-40L}\le e^{-20L}.
\]
Finally, the raw-score comparison and
\[
\frac{\sqrt{mp_xL}+L}{m}
=\sqrt{p_x\delta_{\mathrm p}}+\delta_{\mathrm p}
\]
prove the normalized aggregate bound
\[
\mathbb E_{\mathrm{aux}}\!\left[
(e_x^{\mathrm p})^{j_{\mathrm{mom}}}\mathbf1_{\mathcal G_x^c}
\right]
\le
\left(\frac{H_0}{1024}\right)^{j_{\mathrm{mom}}}
10^{80(j_{\mathrm{mom}}+1)}e^{-20L}
\left(\sqrt{p_x\delta_{\mathrm p}}+\delta_{\mathrm p}\right)^{j_{\mathrm{mom}}}.
\]
Taking \(j_{\mathrm{mom}}=1\) gives the stated aggregate first-moment estimate with numerical constant \(10^{160}\).
% lean: CausalSmith.Stat.OptvalueVanishingoverlapRate.pilotRectangle_bad_probability_le
% lean: CausalSmith.Stat.OptvalueVanishingoverlapRate.pilotRectangle_bad_aggregate_moment_le

The midpoint error satisfies
\[
|b_{ay,x}-q_{ay,x}|\le e_{ay,x}^{\mathrm p},
\qquad
s(b_x)\le p_x+e_x^{\mathrm p}.
\]
Using the clipping-scale square bound from the preceding step, we obtain
\[
\begin{aligned}
\mathbb E_{\mathrm{aux}}
[(S_x^{\mathrm{aw}})^2\mathbf1_{\mathcal G_x^c}]
\le512H_0^2\bigg\{
8e^{-40L}\mathfrak a_x
+\frac{\delta_{\mathrm p}}{\epsilon}
\frac{H_0}{1024}10^{160}e^{-20L}
\left(\sqrt{p_x\delta_{\mathrm p}}+\delta_{\mathrm p}\right)
\bigg\}.
\end{aligned}
\]
Since
\[
\frac{\delta_{\mathrm p}}{\epsilon}
\left(\sqrt{p_x\delta_{\mathrm p}}+\delta_{\mathrm p}\right)
\le2\mathfrak a_x,
\]
where \(\sqrt{p_x\delta_{\mathrm p}}\le p_x+\delta_{\mathrm p}\) and \(\epsilon\le1\), this becomes
\[
\mathbb E_{\mathrm{aux}}
[(S_x^{\mathrm{aw}})^2\mathbf1_{\mathcal G_x^c}]
\le
C_{\mathrm{fail,s}}e^{-20L}\mathfrak a_x,
\]
with
\[
C_{\mathrm{fail,s}}
=
512H_0^2\left(8+2\frac{H_0}{1024}10^{160}\right).
\]

For an occupied cell, \cref{obj:lem:armwise-modulus} gives
\[
|F_\epsilon(b_x)-F_\epsilon(q_x)|
\le
4\sum_{a,y}\frac{p_x}{s_a(q_x)}e_{ay,x}^{\mathrm p}.
\]
The fourth-moment and probability estimates imply
\[
\begin{aligned}
\mathbb E_{\mathrm{aux}}
[(e_{ay,x}^{\mathrm p})^2\mathbf1_{\mathcal G_x^c}]
&\le
\left\{\mathbb E_{\mathrm{aux}}(e_{ay,x}^{\mathrm p})^4\right\}^{1/2}
\Pr(\mathcal G_x^c)^{1/2}\\
&\le
\sqrt{8\cdot10^{20}}
\left(\frac{H_0}{1024}\right)^2
e^{-20L}
\left(\sqrt{q_{ay,x}\delta_{\mathrm p}}+\delta_{\mathrm p}\right)^2.
\end{aligned}
\]
For each arm, the overlap floor gives
\[
\left(\frac{p_x}{s_a(q_x)}\right)^2
\sum_y
\left(\sqrt{q_{ay,x}\delta_{\mathrm p}}+\delta_{\mathrm p}\right)^2
\le8\mathfrak a_x.
\]
Indeed, the left side is at most
\[
2\frac{p_x^2\delta_{\mathrm p}}{s_a(q_x)}
+
4\frac{p_x^2\delta_{\mathrm p}^2}{s_a(q_x)^2}
\le
2\frac{p_x\delta_{\mathrm p}}{\epsilon}
+
4\frac{\delta_{\mathrm p}^2}{\epsilon^2}.
\]
Squaring the four-coordinate sum consequently yields
\[
\mathbb E_{\mathrm{aux}}
[(F_\epsilon(b_x)-F_\epsilon(q_x))^2\mathbf1_{\mathcal G_x^c}]
\le C_{\mathrm{fail,f}}e^{-20L}\mathfrak a_x,
\]
where
\[
C_{\mathrm{fail,f}}
=
1024\sqrt{8\cdot10^{20}}
\left(\frac{H_0}{1024}\right)^2.
\]

Clipping supplies the pathwise bound
\[
|T_x-F_\epsilon(q_x)|
\le
d^{1/4}S_x^{\mathrm{aw}}
+
|F_\epsilon(b_x)-F_\epsilon(q_x)|.
\]
Hence, with
\[
C_{\mathrm d}=2(C_{\mathrm{fail,s}}+C_{\mathrm{fail,f}}),
\]
we have
\[
\mathbb E_{\mathrm{aux}}
[(T_x-F_\epsilon(q_x))^2\mathbf1_{\mathcal G_x^c}]
\le C_{\mathrm d}d^{-4}\mathfrak a_x.
\]
Here
\[
d^{1/2}e^{-20L}\le d^{-4},
\qquad e^{-20L}\le d^{-4}
\]
for every \(d\ge1\). For null cells, \(\mathcal G_x^c\) has probability zero, so the same bound follows directly.
% lean: CausalSmith.Stat.OptvalueVanishingoverlapRate.markedPoissonCellValue_bad_sq_le

\textit{9. Independent-cell risk assembly.}
Define, on the whole auxiliary sample space,
\[
E_x^{\mathrm{good}}
=
(T_x-F_\epsilon(q_x))\mathbf1_{\mathcal G_x},
\qquad
E_x^{\mathrm{bad}}
=
(T_x-F_\epsilon(q_x))\mathbf1_{\mathcal G_x^c}.
\]
Then
\[
T_x-F_\epsilon(q_x)=E_x^{\mathrm{good}}+E_x^{\mathrm{bad}},
\qquad
E_x^{\mathrm{good}}E_x^{\mathrm{bad}}=0.
\]
All these variables have finite second moments by the preceding estimates. Independence of the complete cell statistics gives
\[
\mathbb E_{\mathrm{aux}}
\left(\sum_xT_x-\Psi(\mathbb P)\right)^2
=
\sum_x\operatorname{Var}_{\mathrm{aux}}(T_x)
+
\left\{\sum_x
\mathbb E_{\mathrm{aux}}(T_x-F_\epsilon(q_x))\right\}^2.
\]
The disjoint error decomposition implies
\[
\operatorname{Var}_{\mathrm{aux}}(T_x)
\le
\mathbb E_{\mathrm{aux}}(E_x^{\mathrm{good}})^2
+
\mathbb E_{\mathrm{aux}}(E_x^{\mathrm{bad}})^2.
\]
Applying \((v+w)^2\le2v^2+2w^2\) to the total bias, and Cauchy--Schwarz to its failed-pilot part, therefore gives
\[
\begin{aligned}
\mathbb E_{\mathrm{aux}}
\left(\sum_xT_x-\Psi(\mathbb P)\right)^2
\le{}&
\sum_x\mathbb E_{\mathrm{aux}}(E_x^{\mathrm{good}})^2
+
2\left\{\sum_x\mathbb E_{\mathrm{aux}}E_x^{\mathrm{good}}\right\}^2\\
&+
(1+2d)\sum_x\mathbb E_{\mathrm{aux}}(E_x^{\mathrm{bad}})^2.
\end{aligned}
\]
% lean: CausalSmith.Stat.OptvalueVanishingoverlapRate.independentCell_good_bad_sqRisk_le

\textit{10. Absorption of the pilot scales.}
Assume the active-branch condition \(d/L\le n\epsilon\), and define
\[
\mathfrak r_{\mathrm{aux}}=\frac{d}{m\epsilon L},
\qquad
\mathfrak m_{\mathrm{sum}}
=\sum_x\mathfrak a_x
=\frac L{m\epsilon}+\frac{dL^2}{m^2\epsilon^2}.
\]
Then
\[
0<\mathfrak r_{\mathrm{aux}}\le8,
\qquad
\mathfrak m_{\mathrm{sum}}
\le\frac{L^2+8L^4}{m\epsilon L}.
\]
The second inequality follows by substituting
\(d\le8m\epsilon L\) into its second summand.

For a positive real \(\delta_{\mathrm{dec}}\), define the numerical envelope
\[
\mathcal M(\delta_{\mathrm{dec}})
=
\left(\frac{e^{\delta_{\mathrm{dec}}/2}}{\delta_{\mathrm{dec}}/2}\right)^2
+
8\left(\frac{e^{\delta_{\mathrm{dec}}/4}}{\delta_{\mathrm{dec}}/4}\right)^4.
\]
The elementary inequality \(\log z\le z^u/u\), \(z\ge1\), \(u>0\), applied with \(z=ed\), gives
\[
(L^2+8L^4)d^{-\delta_{\mathrm{dec}}}
\le\mathcal M(\delta_{\mathrm{dec}}).
\]
Set
\[
C_{\mathrm{var}}=\mathcal M(15/16),\qquad
C_{\mathrm{bias}}=\mathcal M(3/8),\qquad
C_{\mathrm{bad,sum}}=\mathcal M(1).
\]
Thus
\[
d^{1/16}\sum_x\mathbb E_{\mathrm{aux}}(S_x^{\mathrm{aw}})^2
\le C_{\mathrm s}C_{\mathrm{var}}\mathfrak r_{\mathrm{aux}}.
\]

The probability-mass identity and Cauchy--Schwarz give
\[
\sum_x\eta_x\le\sqrt{\mathfrak r_{\mathrm{aux}}}+\mathfrak r_{\mathrm{aux}},
\qquad
\sum_x\eta_x^2\le\left(\sum_x\eta_x\right)^2
\le18\mathfrak r_{\mathrm{aux}}.
\]
Also,
\[
\left(\sum_x\mathbb E_{\mathrm{aux}}S_x^{\mathrm{aw}}\right)^2
\le d\sum_x\mathbb E_{\mathrm{aux}}(S_x^{\mathrm{aw}})^2,
\]
and hence
\[
d^{-3/8}
\left(\sum_x\mathbb E_{\mathrm{aux}}S_x^{\mathrm{aw}}\right)^2
\le C_{\mathrm s}C_{\mathrm{bias}}\mathfrak r_{\mathrm{aux}}.
\]
The successful-pilot estimates now imply
\[
\sum_x\mathbb E_{\mathrm{aux}}(E_x^{\mathrm{good}})^2
\le
(C_{\mathrm v}C_{\mathrm s}C_{\mathrm{var}}+18C_{\mathrm e})
\mathfrak r_{\mathrm{aux}},
\]
\[
\left(\sum_x\mathbb E_{\mathrm{aux}}E_x^{\mathrm{good}}\right)^2
\le
2C_{\mathrm b}^2(18+C_{\mathrm s}C_{\mathrm{bias}})
\mathfrak r_{\mathrm{aux}}.
\]
For the failed-pilot term, \(d\,d^{-4}\le1\) gives
\[
d\sum_x\mathbb E_{\mathrm{aux}}(E_x^{\mathrm{bad}})^2
\le C_{\mathrm d}\mathfrak m_{\mathrm{sum}}
\le C_{\mathrm d}C_{\mathrm{bad,sum}}\mathfrak r_{\mathrm{aux}},
\]
so, since \(d\ge1\),
\[
(1+2d)\sum_x\mathbb E_{\mathrm{aux}}(E_x^{\mathrm{bad}})^2
\le3C_{\mathrm d}C_{\mathrm{bad,sum}}\mathfrak r_{\mathrm{aux}}.
\]

Define
\[
\begin{aligned}
C_{\mathrm P}={}&
C_{\mathrm v}C_{\mathrm s}C_{\mathrm{var}}
+4C_{\mathrm b}^2(18+C_{\mathrm s}C_{\mathrm{bias}})
+3C_{\mathrm d}C_{\mathrm{bad,sum}}
+18C_{\mathrm e}+1.
\end{aligned}
\]
This is a positive finite universal constant. The independent-cell assembly yields
\[
\mathbb E_{\mathrm{aux}}
\left(\sum_xT_x-\Psi(\mathbb P)\right)^2
\le C_{\mathrm P}\frac{d}{m\epsilon L}.
\]
Finally, define the projected auxiliary statistic by
\[
V_{\mathrm{aux}}=\Pi_{[0,1]}\left(\sum_xT_x\right).
\]
Since \(\Psi(\mathbb P)\in[0,1]\), projection contracts its squared error, giving
\[
\mathbb E_{\mathrm{aux}}
(V_{\mathrm{aux}}-\Psi(\mathbb P))^2
\le C_{\mathrm P}\frac{d}{m\epsilon L}.
\]
% lean: CausalSmith.Stat.OptvalueVanishingoverlapRate.activeBranch_good_bad_approximation_sqRisk_rate_le
% lean: CausalSmith.Stat.OptvalueVanishingoverlapRate.markedPoissonStatistic_goodPilotEstimates_rate_le
% lean: CausalSmith.Stat.OptvalueVanishingoverlapRate.markedPoissonStatistic_active_rate_tuning_exists

\textit{11. Fixed-sample transfer.}
Set
\[
D_0=2,\qquad
C_\star=3+8\log(2e)+8C_{\mathrm P}.
\]
Then \(D_0\ge2\) and \(0<C_\star<\infty\). For the admissible theorem parameters \(n\ge1,d\ge2\), the finite sample alphabet makes the estimator in \cref{obj:def:armwise-estimator} measurable. Its auxiliary variables are integrated by the specified finite sum, so it is a deterministic function of the observations.

Fix an admissible observed law and its sampling law from \cref{obj:ass:iid-sampling}. First consider \(d<D_0\). With the chosen cutoff this branch has an empty parameter range, because \(d\ge2=D_0\).

Next suppose \(d\ge D_0\) and \(n\epsilon<d/L\). The estimator equals \(1/2\), and
\[
\mathbb E_{\mathbb P}
\left(\widehat V_{n,d,\epsilon}^{\mathrm{AF}}-\Psi(\mathbb P)\right)^2
\le\frac14,
\qquad
\min\left\{1,\frac{d}{n\epsilon L}\right\}=1.
\]
This proves the desired bound on that branch.

On the active branch, put the auxiliary count and marks alongside the fixed observed sample. Conditional averaging of squared error gives
\[
\begin{aligned}
\mathbb E_{\mathbb P}
\left(\widehat V_{n,d,\epsilon}^{\mathrm{AF}}-\Psi(\mathbb P)\right)^2
\le{}&
\mathbb E_{\mathrm{aux}}
\left[
(V_{\mathrm{aux}}-\Psi(\mathbb P))^2
\mathbf1_{\{\mathsf M\le n\}}
\right]\\
&+\Pr(\mathsf M>n)\Psi(\mathbb P)^2.
\end{aligned}
\]
Indeed, the auxiliary mixture has value zero on overflow. For each \(k\le n\), the first \(k\) observations have law \(\mathbb P^{\otimes k}\); consequently the integrated nonoverflow loss is exactly the corresponding part of the unrestricted finite-Poisson experiment.

The quarter-mean Poisson cap satisfies
\[
\Pr(\mathsf M>n)\le e^{-n/2}.
\]
For example, exponential Markov inequality at \(t=\log4\) gives
\[
\Pr(\mathsf M>n)
\le e^{-n(\log4-3/4)}
\le e^{-n/2}.
\]
Also, \(e^{n/2}\ge n/2\), \(L\le d\), and \(\epsilon\le1/2\) imply
\[
e^{-n/2}\le\frac2n
\le2\frac{d}{n\epsilon L}.
\]
Since \(\Psi(\mathbb P)^2\le1\), the auxiliary risk bound and \(m=n/8\) therefore give
\[
\mathbb E_{\mathbb P}
\left(\widehat V_{n,d,\epsilon}^{\mathrm{AF}}-\Psi(\mathbb P)\right)^2
\le
(8C_{\mathrm P}+2)\frac{d}{n\epsilon L}.
\]
The estimator and target both lie in \([0,1]\), so the risk is also at most one. Splitting according as \(d/(n\epsilon L)\ge1\) or \(d/(n\epsilon L)\le1\), we obtain
\[
\mathbb E_{\mathbb P}
\left(\widehat V_{n,d,\epsilon}^{\mathrm{AF}}-\Psi(\mathbb P)\right)^2
\le
(8C_{\mathrm P}+2)
\min\left\{1,\frac{d}{n\epsilon L}\right\}
\le
C_\star\min\left\{1,\frac{d}{n\epsilon L}\right\}.
\]
% lean: CausalSmith.Stat.OptvalueVanishingoverlapRate.armwiseEstimator_active_le_markedPoisson_sqRisk
% lean: CausalSmith.Stat.OptvalueVanishingoverlapRate.armwiseEstimator_allBranches_rate_le

\textit{12. Uniformity over sampling laws.}
The preceding estimates hold for every admissible observed law and every family of sampling laws satisfying \cref{obj:ass:iid-sampling}. If the admissible class is empty, its worst-risk value, with the empty-class convention zero, satisfies the bound because the right side is nonnegative. Otherwise, taking the supremum of the pointwise bound gives
\[
\sup_{\mathbb P\in\mathcal D_{d,\epsilon}^{\mathrm{obs}}}
\mathbb E_{\mathbb P}
\left[
\left\{
\widehat V_{n,d,\epsilon}^{\mathrm{AF}}-\Psi(\mathbb P)
\right\}^{2}
\right]
\le
C_\star\min\left\{1,\frac{d}{n\epsilon L_d}\right\}.
\]
Together with measurability and determinism, this proves the claim.
% lean: CausalSmith.Stat.OptvalueVanishingoverlapRate.armwiseEstimator_allBranches_worstRisk_le
% lean: CausalSmith.Stat.OptvalueVanishingoverlapRate.observable_upper
\end{proof}
